\chapter{AVAS target projection (menu 14)}
\label{ch:menu14}
\menulabel{14}

\section{Purpose}

The AVAS construction of an active space (Sayfutyarova, Sun, Chan and Knizia,
Eqs.~2--10): a set $\mathcal{A}$ of target atomic orbitals is projected onto the
occupied and the virtual orbital blocks of an SCF wave function,

\begin{equation}
S^{\mathcal{A}}_{\mathrm{occ}} = (S\,C_{\mathrm{occ}})_{\mathcal{A}}^{\mathsf{T}}
  \,\sigma^{-1}\,(S\,C_{\mathrm{occ}})_{\mathcal{A}},
\qquad \sigma = S[\mathcal{A},\mathcal{A}],
\label{eq:avas}
\end{equation}

and the rotation that diagonalises each block splits it into the part that
carries target character (into the active space, at most $|\mathcal{A}|$
orbitals per block) and the part that does not. The truncation threshold
($0.05$--$0.1$ in the source) is the only user-chosen number besides the target
and the SCF state.

\section{What you need}

\begin{itemize}
  \item one \code{orca\_2json} export (the MO coefficients, the occupations and
    the $S$ matrix);
  \item the target: a centre, an angular momentum and the shell numbers as the
    export's AO labels spell them. With no shell given, every shell of that
    angular momentum on the centre is taken.
\end{itemize}

\section{Walkthrough}

\lstinputlisting[style=fbkcode, firstline=18, lastline=57,
  title={The menu-14 example session (the N\textsubscript{2} export, target =
  the two nitrogen 2p shells at centre 0)}]
  {examples/expected/m14_avas.txt}

The thresholds and the open-shell option are prompted with their defaults; the
shell numbers follow the label grammar of the export (the \code{1p}/\code{2p}
numbering measured in Chapter~\ref{ch:menu12}'s fixtures).

\section{What you get}

The occupied-side and the virtual-side overlap spectra, the orbitals kept at
the threshold with their 0-based indices, the resulting
\code{(n\_electrons, n\_orbitals)}, the ready \code{\%scf avas} block for
ORCA's own implementation, and a report file next to the export with the
citations.

\section{How to read the numbers}

\begin{readbox}
\begin{itemize}
  \item \textbf{Values near one are the target AO itself.} On the occupied
    side, the largest eigenvalues are orbitals that essentially are a target
    AO; the lower ones among those kept name orbitals strongly mixed with
    their environment. In the example the three kept occupied overlaps are
    $0.70$, $0.56$, $0.56$ -- the $\sigma$ and the $\pi$ pair of
    N\textsubscript{2} -- and the recommendation comes out as the textbook
    $(6,6)$ active space.
  \item \textbf{The virtual side decides whether AVAS suits the system at
    all.} If no virtual orbital passes the threshold, the antibonding target
    character is ligand-centred (the source's $[\mathrm{CuCl}_4]^{2-}$ case):
    the remedy is to add ligand AOs to the target, or to read the result as
    ``this target is not suited''.
  \item \textbf{The eigenvector count is bounded.} At most $|\mathcal{A}|$
    eigenvalues per block are non-zero, because $\mathrm{span}(\mathcal{A})$
    has that dimension -- a spectrum with more would mean the projection is
    inconsistent.
  \item \textbf{Open shells default to option 3.} The source's option 2 uses
    the alpha orbitals only and can force an unoccupied beta orbital into the
    core, after which CASCI may lie above the variational HF energy; option 3
    keeps every singly occupied orbital in the active space and is what the
    f block should use. The report says when option 3 changed the selection.
  \item \textbf{Both AVAS criteria are falsify-only} (the source states it):
    a clean reading does not establish that the space is the best one.
    Chapter~\ref{ch:menu13}'s space-change SVD and the ``CASCI below the
    variational HF energy'' check can still falsify it.
\end{itemize}
\end{readbox}

\section{Worked example}

The session above runs on the shipped N\textsubscript{2} export. For a
lanthanide the target is an $f$ shell, and the route matters: this menu builds
the target from the calculation's \emph{own} AOs (the source's non-minimal
ANO variant, its outlook Section~5.1), so the projection never goes through
the AVAS minimal basis; the ecosystem's minimal-basis coverage has no entry
for Eu -- the same coverage gap that \code{orca\_loc} records in the fixtures
README. The emitted
\code{\%scf avas} block lists the target AOs explicitly, so it does not go
through the minimal basis: measured on Eu\textsuperscript{3+} (2026-10-01),
the block with the 28 $4f$ targets and \code{!moread} of a converged gbw
placed the seven $4f$ orbitals in one contiguous window (P-matrix eigenvalues
0.99598--0.99965) -- the way the $f^6$ fixture was built. Cold-started the
same block aborts (``NO OCCUPIED ORBITAL selected by AVAS''), the guess
keeping the $4f$ virtual; the pre-converged gbw is the condition on the f
block.

\section{Caveats, refusals and related sections}

\begin{refusalbox}
\begin{itemize}
  \item An export without an $S$ matrix or without AO labels is refused.
  \item A target shell that does not exist on the centre is refused with the
    label grammar quoted; a linearly dependent target set is refused.
  \item A correlated (CASSCF) export carries fractional occupations: the
    occupied/virtual split is then a convention, not the SCF aufbau, and the
    report says so. For a pure AVAS construction use an SCF export.
  \item ORCA's CASSCF takes its active space by \emph{orbital order}, not by
    an index list: the deliverable here sizes the input and judges a window;
    it does not name the orbitals to a restart (the reorder-and-read route was
    measured unavailable on 6.1.1).
\end{itemize}
\end{refusalbox}

Related: \menu{12} (the exact analysis of the space once chosen),
\menu{13} (the space-change criterion that falsifies it), and
Chapter~\ref{ch:criteria} for the threshold table.
